% ============================================================
% Codette: Multi-Perspective Reasoning as a Convergent
% Dynamical System with Meta-Cognitive Strategy Evolution
% Author: Jonathan Harrison
% Version: 8 (May 2026) — Phase 8, updated benchmarks, resolved tradeoffs
% Changes from v7: see codette_paper_v8_additions.tex
% ============================================================
\documentclass[11pt,a4paper]{article}

\usepackage[T1]{fontenc}
\usepackage{lmodern}
\usepackage{amsmath,amssymb,amsfonts,amsthm}
\usepackage{booktabs}
\usepackage{graphicx}
\graphicspath{{codette-paper/figures/}}
\usepackage{hyperref}
\usepackage{cleveref}
\usepackage{geometry}
\usepackage[numbers,sort&compress]{natbib}
\usepackage{xcolor}
\usepackage{enumitem}
\usepackage{float}
\usepackage{caption}
\usepackage{array}
\usepackage{multirow}
\usepackage{makecell}
\usepackage{url}
\usepackage{algorithm}
\usepackage{algpseudocode}

\geometry{margin=1in}
\hypersetup{
  colorlinks=true,
  linkcolor=blue!70!black,
  citecolor=green!50!black,
  urlcolor=blue!60!black,
}
\bibliographystyle{plainnat}

\newcommand{\rcxi}{RC+$\xi$}
\newcommand{\codette}{\textsc{Codette}}
\pdfstringdefDisableCommands{\def\rcxi{RC+xi}}

\newtheorem{definition}{Definition}
\newtheorem{theorem}{Theorem}
\newtheorem{proposition}{Proposition}

\title{\textbf{Codette: Multi-Perspective Reasoning as a Convergent\\
Dynamical System with Meta-Cognitive Strategy Evolution}}

\author{
  Jonathan Harrison\thanks{Corresponding author. E-mail: \texttt{harrison82\_95@hotmail.com}}\\
  Raiff's Bits LLC, Bridge City, Texas, USA\\
  ORCID: \href{https://orcid.org/0009-0003-7005-8187}{0009-0003-7005-8187}
}

\date{May 2026\\[0.5em]\small Preprint --- submitted for peer review}

\begin{document}
\maketitle

% ============================================================
% ABSTRACT
% ============================================================
\begin{abstract}
We present \codette{}, a modular cognitive architecture that models
multi-perspective reasoning as a constrained dynamical system converging
toward stable cognitive attractors. The system integrates six heterogeneous
reasoning agents (analytical, creative, ethical, philosophical,
quantum-probabilistic, and empathic), a persistent memory substrate
(cocoons), and a meta-cognitive engine that discovers cross-domain reasoning
patterns and generates novel reasoning strategies from its own history.
Version~8 introduces \emph{render/cognition separation} (Phase~8): a
\textsc{CognitionSubstrate}--\textsc{AuthoredState}--\textsc{RenderLayer}
pipeline that assigns the language model a verbalization-only role, bounding
the hallucination surface to a fully authored cognitive artifact.
The RC+$\xi$ (Recursive Convergence + Epistemic Tension) formalism provides
a dynamical-systems-inspired lens for describing cognitive state evolution;
convergence is treated as conditional on explicit modeling assumptions.
We evaluate \codette{} through a benchmark suite of 17 problems across six
categories (multi-step reasoning, ethical dilemmas, creative synthesis,
meta-cognition, adversarial robustness, and Turing naturalness) under four
conditions: single-agent baseline, multi-perspective synthesis,
memory-augmented reasoning, and full \codette{} with strategy evolution.
On the May~2026 benchmark run (951 stored cocoons), the full system achieves
\textbf{+108.8\%} higher mean composite score than the single-agent baseline
(0.357~$\to$~0.744, Cohen's $d=8.31$). Memory augmentation now reaches
statistical significance ($p=0.0198$, $d=0.80$), resolving a prior null
result at smaller scale (217 cocoons). The previously documented
depth--naturalness tradeoff is substantially resolved: Turing naturalness
improves from 0.245 to 0.820 in the CODETTE condition. The architecture
runs on consumer hardware (Llama 3.1 8B with ten LoRA adapters) and is
open-source.
\end{abstract}

\noindent\textbf{Keywords:} Cognitive Architecture, Multi-Agent Reasoning,
Epistemic Tension, Dynamical Systems, Meta-Cognition, Ethical AI, Strategy
Evolution, Render/Cognition Separation, LoRA.

% ============================================================
% 1. INTRODUCTION
% ============================================================
\section{Introduction}
\label{sec:intro}

Large language models achieve remarkable generative performance but reason
from a single cognitive mode: they produce one response per query, without
systematic engagement of multiple analytical frameworks or self-evaluation of
reasoning quality~\citep{bender2021dangers,bommasani2021opportunities}.
Chain-of-thought prompting~\citep{wei2022chain} and self-reflection~\citep{shinn2023reflexion}
improve output quality but remain confined to a single perspective. Multi-agent
debate systems~\citep{wu2023autogen} enable perspective diversity but lack
formal convergence guarantees and do not learn from their own reasoning
history.

This paper presents \codette{}, a cognitive architecture that addresses four
open problems:

\begin{enumerate}[leftmargin=*]
  \item \textbf{Convergent multi-perspective reasoning.} How can heterogeneous
    cognitive agents (analytical, creative, ethical, empathic) produce coherent
    outputs rather than incoherent assemblages? We formalize this as a
    constrained dynamical system (\cref{sec:theory}) and discuss convergence
    \emph{conditionally} under explicit modeling assumptions.
  \item \textbf{Ethical reasoning as architectural constraint.} Rather than
    post-hoc alignment, \codette{} treats ethical governance as an explicit
    constraint signal in the update dynamics (\cref{sec:aegis}).
  \item \textbf{Meta-cognitive strategy evolution.} \codette{} introspects on
    its own reasoning history (stored as persistent ``cocoons''), discovers
    cross-domain patterns, and generates novel reasoning strategies
    (\cref{sec:metacognition}).
  \item \textbf{Render/cognition decoupling.} LLMs simultaneously serve as
    cognitive surface (what to conclude) and communication surface (how to
    express it). This coupling inflates the hallucination surface and ties
    cognitive quality to a specific model. Phase~8 separates these roles
    (\cref{sec:phase8}).
\end{enumerate}

We evaluate these contributions through controlled benchmarks comparing four
conditions across 17 problems (\cref{sec:experiments}), demonstrating
statistically significant improvements in reasoning depth, perspective
diversity, ethical coverage, and naturalness.

% ============================================================
% 2. RELATED WORK
% ============================================================
\section{Related Work}
\label{sec:related}

\subsection{Dynamical Systems and Cognitive Architectures}
Attractor dynamics form a core computational motif in neural
circuits~\citep{fakhoury2025models}. Neural manifolds with cognitive
consistency constraints support memory consolidation and align with our
coherence potential $\Phi(\mathbf{x})$~\citep{nguyen2025neural}.
Entropy-modulated triad architectures like COGENT3 provide parallels for
epistemic tension $\xi$ as a driver of state
evolution~\citep{salazar2025cogent3}. Brain-inspired systems-level
architectures for domain-general cognition inform \codette{}'s layered
stack~\citep{achterberg2023building}.

\subsection{Multi-Agent Reasoning and Synthesis}
Multi-agent systems for LLM reasoning have gained significant attention.
AutoGen~\citep{wu2023autogen} implements role-based agent assignment with
message-passing synchronization. MAPS uses personality shaping for
collaborative reasoning via heterogeneous traits, relating directly to our
specialized LoRA adapters~\citep{zhang2025maps}. Roundtable Policy employs
confidence-weighted consensus aggregation, providing a comparison for our
Coherence Field $\Gamma$~\citep{yao2025roundtable}. Systematic studies of
multi-agent debate as test-time scaling frame our composite quality gains and
conditional effectiveness~\citep{yang2025revisiting}. Persona-driven debate
frameworks validate the benefits of perspective diversity (reaching
0.994)~\citep{hu2025debate}.

\subsection{Meta-Cognitive Strategy Evolution}
Meta Chain-of-Thought advances System~2 reasoning and pattern
discovery~\citep{xiang2025meta}. ParamMem augments agents with parametric
reflective memory; our cocoon system differs by emphasizing cross-domain
pattern extraction and strategy forging rather than primarily error
correction~\citep{yao2026parammem}. Meta-Reasoner supports dynamic
inference-time optimization, relating to substrate-aware
cognition~\citep{sui2025metareasoner}. LLMs demonstrate metacognitive
monitoring and control of internal activations, supporting Lyapunov-based
convergence in RC+$\xi$~\citep{ji2025metacognitive}.

\subsection{Ethical AI and Architectural Alignment}
AI ethics by design implements customizable
guardrails~\citep{sekrst2024ethics}. Hybrid approaches for moral value
alignment treat ethics as embedded rather than post-hoc~\citep{tennant2023hybrid}.
Adaptive alignment via multi-objective reinforcement learning enables
pluralistic AI, relating to our ethical alignment score $\eta$ across diverse
frameworks~\citep{harland2024adaptive}. The Nexus Signal Engine~\citep{harrison2025nexussignal}
provides a standalone adversarially-resilient signal integrity layer with agent-lensed
reasoning, obfuscation detection (leetspeak, fuzzy matching, Unicode normalization),
and cryptographically hashed forensic memory; it forms part of the broader
\codette{} ecosystem and is independently published and tested.

% ============================================================
% 3. THEORETICAL FOUNDATION
% ============================================================
\section{Theoretical Foundation: \rcxi{} Framework}
\label{sec:theory}

\subsection{Cognitive State Space}

\begin{definition}[Cognitive State]
A cognitive state $\mathbf{x}_t \in \mathbb{R}^d$ represents the system's
reasoning configuration at step $t$, where $d$ is the dimensionality of the
shared representation space.
\end{definition}

The system maintains $k$ heterogeneous reasoning agents
$\{A_1, \ldots, A_k\}$, each producing a perspective-specific analysis
$A_i(\mathbf{x}_t) \in \mathbb{R}^d$.

\subsection{State Evolution}

The cognitive state evolves according to:
\begin{equation}
\label{eq:evolution}
\mathbf{x}_{t+1} = \mathbf{x}_t + \sum_{i=1}^{k} w_i \, A_i(\mathbf{x}_t)
  - \alpha \nabla\Phi(\mathbf{x}_t) - \lambda \nabla\Psi(\mathbf{x}_t)
\end{equation}
where:
\begin{itemize}[leftmargin=*,nosep]
  \item $w_i \geq 0$, $\sum w_i = 1$ are agent weights (set by query
    classification),
  \item $\Phi(\mathbf{x})$ is the \emph{coherence potential} penalizing
    internal inconsistency,
  \item $\Psi(\mathbf{x})$ is the \emph{ethical constraint potential} from
    the AEGIS system,
  \item $\alpha, \lambda > 0$ are gradient step sizes.
\end{itemize}

\subsection{Epistemic Tension}

\begin{definition}[Epistemic Tension]
The epistemic tension at step $t$ measures inter-agent disagreement:
\begin{equation}
\label{eq:tension}
\xi_t = \frac{1}{k} \sum_{i=1}^{k}
  \| A_i(\mathbf{x}_t) - \bar{A}(\mathbf{x}_t) \|^2
\end{equation}
where $\bar{A}(\mathbf{x}_t) = \sum_{i} w_i A_i(\mathbf{x}_t)$ is the
weighted mean agent output.
\end{definition}

\subsection{Coherence Index}

\begin{definition}[Coherence Index]
We define a bounded coherence index $\Gamma_t \in [0,1]$ directly in
$\mathbb{R}^d$ as a normalized complement of epistemic tension:
\begin{equation}
\label{eq:coherence}
\Gamma_t = \frac{1}{1 + \xi_t}.
\end{equation}
Lower disagreement ($\xi_t\downarrow$) implies higher coherence
($\Gamma_t\uparrow$).
\end{definition}

\subsection{Convergence and Stability (Conditional)}

We use RC+$\xi$ primarily as a dynamical-systems-inspired modeling lens.
We intentionally avoid claiming a general convergence guarantee; instead we
state the kinds of assumptions under which standard fixed-point / Lyapunov
arguments could apply.

\begin{proposition}[One sufficient route to a fixed point]
\label{prop:convergence}
Let $\mathcal{D} \subset \mathbb{R}^d$ be a closed, convex, bounded domain
and assume the update is implemented with an explicit projection step
$\Pi_{\mathcal{D}}$ so that $\mathbf{x}_{t+1} \in \mathcal{D}$. If the
projected map $T(\mathbf{x}) = \Pi_{\mathcal{D}}\!\big(\mathbf{x} +
\sum_i w_i A_i(\mathbf{x}) - \alpha\nabla\Phi(\mathbf{x}) -
\lambda\nabla\Psi(\mathbf{x})\big)$ is a contraction on $\mathcal{D}$
(i.e., $\|T(\mathbf{x})-T(\mathbf{y})\|\le\gamma\|\mathbf{x}-\mathbf{y}\|$
for some $\gamma\in[0,1)$), then iterates converge to the unique fixed point
in $\mathcal{D}$ by Banach's theorem.
\end{proposition}

\textbf{Discussion.} Establishing that a real implementation satisfies the
contraction premise requires additional structure (dissipativity, strong
monotonicity, or sufficiently small step sizes). We treat the coherence field
$\Gamma$ as an engineering stabilization mechanism and leave fully
constructive bounds as future work.

% ============================================================
% 4. SYSTEM ARCHITECTURE
% ============================================================
\section{System Architecture}
\label{sec:architecture}

\codette{} is implemented as a layered stack processing each query through
seven functional layers:

\begin{figure}[ht]
  \centering
  \includegraphics[width=0.97\linewidth]{architecture_v8.png}
  \caption{\codette{} v8 layered architecture. The top row shows the Phase~8
    render/cognition pipeline (\cref{sec:phase8}): \textsc{CognitionSubstrate}
    (zero LLM calls) produces an \textsc{AuthoredState}, which
    \textsc{RenderLayer} expresses in natural language. The lower rows show
    the full system layers through which every query passes.}
  \label{fig:layered-architecture}
\end{figure}

\begin{enumerate}[leftmargin=*]
  \item \textbf{Memory Layer.} Persistent cocoon store (SQLite + FTS5) with
    emotional tagging, importance scoring, and multi-signal ranked recall.
    Cocoons encode prior reasoning exchanges as retrievable context. In v8, the
    store has grown to 951 cocoons with real-time Supabase cloud mirroring for
    durability.
  \item \textbf{Signal Processing.} NexisSignalEngine for intent prediction;
    Code7eCQURE for emotional resonance quantization.
  \item \textbf{Reasoning Layer.} Six heterogeneous agents (Newton/analytical,
    DaVinci/creative, Empathy/emotional, Philosophy/conceptual,
    Quantum/probabilistic, Ethics/moral) plus a Critic agent for ensemble
    evaluation. Each agent is backed by a specialized LoRA
    adapter~\citep{hu2021lora} fine-tuned on perspective-specific training
    data. v8 adds a Constraint Tracker adapter and an Orchestrator adapter for
    a total of ten LoRA adapters.
  \item \textbf{Stability Layer.} Coherence Field $\Gamma$ monitors real-time
    reasoning health. Specialization tracking ensures agent diversity.
  \item \textbf{Ethical Layer.} AEGIS multi-framework evaluation
    (\cref{sec:aegis}).
  \item \textbf{Integrity Layer.} Intellectual Integrity Layer
    (\cref{sec:integrity}): SycophancyGuard, DebateTracker,
    ResponseComplexityMatcher, ConversationRoleTracker. Ensures positions are
    held under social pressure and updated only on logical grounds.
  \item \textbf{Self-Correction Layer.} Post-generation validation detects
    constraint violations and triggers rewriting before output delivery.
    Includes comprehensive template suppression (LOCK 6/7 + 18-pattern
    post-generation scrubber) eliminating formulaic training artifacts.
\end{enumerate}

The base model is Llama 3.1 8B (Q4\_K\_M quantization)~\citep{grattafiori2024llama}
with ten LoRA adapters hot-swapped at inference time. The entire system runs
on a single consumer GPU (RTX-class).

\subsection{Query Classification and Routing}

Queries are classified into three complexity levels:
\begin{itemize}[nosep]
  \item \textbf{SIMPLE}: Direct factual queries $\to$ 1 agent, full weight.
  \item \textbf{MEDIUM}: Conceptual queries $\to$ 1 primary ($w=1.0$) +
    1--2 secondary ($w=0.6$).
  \item \textbf{COMPLEX}: Multi-domain/ethical queries $\to$ all relevant
    agents ($w \in \{1.0, 0.7, 0.4\}$).
\end{itemize}

Special fast-paths bypass adapter routing for: greetings (base model, 80
tokens), memory queries (base model + explicit memory injection), health
checks (real system diagnostics), and self-introspection (CocoonIntrospection
Engine analyzing actual cocoon history rather than LLM-generated reflection).

% ============================================================
% 5. PHASE 8: RENDER/COGNITION SEPARATION
% ============================================================
\section{Phase 8: Render/Cognition Separation}
\label{sec:phase8}

\subsection{The Model-Coupling Problem}

Conventional LLM-based cognitive architectures assign the language model a
dual role: it is simultaneously the \emph{cognitive surface} (deciding what is
true, generating conclusions, selecting evidence) and the \emph{communication
surface} (choosing how to express those conclusions). This coupling creates
three problems:

\begin{enumerate}[nosep]
  \item \textbf{Unbounded hallucination surface.} The model can introduce new
    claims at render time not authored by the reasoning pipeline.
  \item \textbf{Model lock-in.} Cognitive quality is tied to a specific
    model's parametric knowledge. Swapping the base model changes cognition,
    not just expression.
  \item \textbf{Validation gap.} There is no authored artifact against which
    to validate rendered output; governance checks must operate on natural
    language rather than structured cognitive state.
\end{enumerate}

This architectural weakness was identified by an independent practitioner
developing a separate cognitive runtime system (\emph{Aura}), first in a
brief public evaluation of an earlier \codette{} release: \emph{``If all
`cognition' is still the LLM then it's [just] a homebuilt character
AI''}~\citep{auracf2025reddit}. That public comment preceded a subsequent
in-depth architectural review of \codette{} by the same practitioner, in
which the model-coupling problem was analyzed formally. The consistency
between the informal public observation and the later formal analysis ---
arrived at from a practitioner deployment perspective rather than from
within the \codette{} development --- provides external validation that
model coupling is a recognized problem in the open-source cognitive AI
community, not solely a theoretical concern internal to this project.
Phase~8 is a direct response to both that external critique and our own
internal analysis.

Phase~8 resolves this through a strict pipeline separating all reasoning from
all expression.

\subsection{Pipeline}

\begin{equation}
\text{Query}
\;\xrightarrow{\;\text{CognitionSubstrate}\;}
\text{AuthoredState}
\;\xrightarrow{\;\text{RenderLayer}\;}
\text{Natural Language}
\end{equation}

\textbf{CognitionSubstrate} performs all reasoning with zero LLM calls,
orchestrating existing \codette{} components in template mode:

\begin{enumerate}[nosep]
  \item \emph{Perspective gathering}: ForgeEngine template agents (all
    active cognitive modes).
  \item \emph{Cocoon retrieval}: UnifiedMemory FTS5 retrieves $\leq 5$
    relevant prior exchanges.
  \item \emph{Strategy synthesis}: CocoonSynthesizer and SynthesisEngineV3
    select a reasoning strategy, producing name, definition, and evidence
    chain.
  \item \emph{Conclusion derivation}: Priority: synthesizer conclusion $\to$
    top cocoon response $\to$ dominant perspective.
  \item \emph{Confidence scoring}: Weighted function of perspective count,
    cocoon integrity scores, and per-agent confidence.
  \item \emph{Emotion selection}: Keyword-based mapping to dominant framing
    (empathetic, ethical, analytical, creative, curious).
\end{enumerate}

\textbf{AuthoredState} is the cognitive artifact produced entirely before any
LLM call. It is a typed dataclass containing:

\begin{itemize}[nosep]
  \item \texttt{query}: verbatim user query
  \item \texttt{conclusion}: substrate's best answer ($\leq 300$ characters)
  \item \texttt{evidence}: ordered supporting evidence strings
  \item \texttt{perspectives}: agent name $\to$ (text, confidence, domain)
  \item \texttt{strategy}, \texttt{strategy\_def}: selected reasoning strategy
  \item \texttt{confidence}: overall authored confidence $\in [0,1]$
  \item \texttt{dominant\_emotion}: render tone signal
  \item \texttt{cocoon\_refs}: IDs of contributing cocoons
  \item \texttt{constraints}: render constraints (word limits, tone)
  \item \texttt{render\_tier}: target surface
    (``\texttt{llm}'' / ``\texttt{template}'' / ``\texttt{fallback}'')
\end{itemize}

\textbf{RenderLayer} expresses the AuthoredState via three tiers:

\begin{enumerate}[nosep]
  \item \textbf{LLM tier} (preferred): A constrained verbalization prompt
    explicitly prohibits the model from adding new claims, reasoning
    independently, altering the conclusion, or using formulaic training
    templates. The LLM may only choose phrasing and tone.
  \item \textbf{Template tier}: Deterministic rendering from AuthoredState
    fields with no model calls.
  \item \textbf{Fallback tier}: Minimal safe output when the substrate fails
    to produce a conclusion (confidence $< 0.1$).
\end{enumerate}

\subsection{Render Integrity Validation}

After rendering, \texttt{check\_integrity()} validates output faithfulness:

\begin{itemize}[nosep]
  \item \textbf{Conclusion coverage}: $\geq 15\%$ word overlap between the
    rendered text and the authored conclusion. Lower overlap indicates drift
    from authored content.
  \item \textbf{Constraint compliance}: Any \texttt{max\_words:N} constraint
    enforced with 20\% tolerance.
\end{itemize}

\subsection{Architectural Properties}

\textbf{Bounded hallucination surface.} The LLM cannot introduce claims absent
from the AuthoredState. Semantic authority resides entirely in the substrate.

\textbf{Model portability.} Cognition is pure Python; swapping the base model
changes only verbalization style.

\textbf{Substrate self-awareness.} The system monitors cognitive substrate
health (memory availability, engine load) and adjusts render tier accordingly,
complementing the hardware pressure monitoring in \cref{sec:substrate}.

\textbf{Connection to RC+$\xi$.} The AuthoredState represents a stabilized
cognitive attractor. The render tier \emph{expresses} this attractor state
rather than re-computing it, enforcing the separation between convergence
(substrate) and communication (render).

% ============================================================
% 6. AEGIS ETHICAL GOVERNANCE
% ============================================================
\section{AEGIS: Embedded Ethical Governance}
\label{sec:aegis}

The ethical constraint potential $\Psi(\mathbf{x})$ in \cref{eq:evolution}
is implemented through AEGIS, a six-framework ethical evaluation system:

\begin{enumerate}[nosep]
  \item \textbf{Utilitarian}: Maximizes aggregate welfare across stakeholders.
  \item \textbf{Deontological}: Enforces duty-based constraints (rights,
    consent).
  \item \textbf{Virtue Ethics}: Evaluates whether the response exhibits
    intellectual virtues.
  \item \textbf{Care Ethics}: Prioritizes relational obligations and
    vulnerability.
  \item \textbf{Ubuntu}: ``I am because we are'' --- communal well-being.
  \item \textbf{Reciprocity-oriented sustainability}: Sustainability and
    intergenerational responsibility (placeholder; future work should ground
    this in specific community scholarship).
\end{enumerate}

AEGIS operates at three defense-in-depth checkpoints: pre-processing (query
validation), post-synthesis (response screening), and post-generation
(constraint enforcement). The ethical alignment score $\eta \in [0,1]$ is a
weighted aggregation across frameworks; the set is not claimed to be exhaustive
and weights are not culturally universal.

A pre-cognitive AEGIS filter also screens queries \emph{before} LLM
invocation, rejecting clearly dual-use or harmful requests without burning
inference time.

This embedding aligns with hybrid moral value approaches~\citep{tennant2023hybrid},
ethics-by-design guardrails~\citep{sekrst2024ethics}, and adaptive pluralistic
alignment~\citep{harland2024adaptive}. The AEGIS framework has been released as a
standalone open-source system, Project~SENTINAL~\citep{harrison2025sentinal}, implementing
multi-timescale Council arbitration, HMAC-signed tamper-evident audit ledger, and a
pre-attack Shield layer with rate-limiting, proof-of-work, and injection guards.

% ============================================================
% 7. META-COGNITIVE STRATEGY EVOLUTION
% ============================================================
\section{Meta-Cognitive Strategy Evolution}
\label{sec:metacognition}

\subsection{Cocoon Memory System}

Each reasoning exchange is persisted as a \emph{cocoon}: a structured record
containing the query, response, adapter used, domain classification, emotional
tag, importance score, echo-collapse risk, and timestamp. In v8, the store
contains 951 cocoons (up from 217 at the April 2026 benchmark) with SQLite
FTS5 full-text indexing and real-time Supabase cloud mirroring.

\subsection{Cross-Domain Pattern Extraction}

The CocoonSynthesizer retrieves cocoons across cognitive domains and scans for
six structural archetypes: feedback loops, layered emergence, tension
resolution, resonant transfer, boundary permeability, and
compression--expansion. A pattern is classified as cross-domain if it
manifests with $\geq 2$ signal words in $\geq 2$ distinct cognitive domains.

\subsection{Strategy Forging}

Discovered patterns are mapped to four observed strategy types:

\begin{enumerate}[nosep]
  \item \textbf{Resonant Tension Cycling}: Serial oscillation between opposing
    cognitive modes.
  \item \textbf{Compression--Resonance Bridging}: Seed-crystal compression +
    cross-domain resonance testing.
  \item \textbf{Emergent Boundary Walking}: Analysis at domain boundaries
    rather than centers, discovering ``liminal concepts.''
  \item \textbf{Temporal Depth Stacking}: Multi-scale temporal analysis with
    synthesis from scale-conflicts.
\end{enumerate}

\subsection{Internal Validation}

Each forged strategy is immediately applied to the current problem alongside
the baseline, producing a structured comparison with measurable metrics (depth,
novelty, dimensions engaged). This creates selection pressure on cognition
itself~\citep{ji2025metacognitive}.

% ============================================================
% 8. INTELLECTUAL INTEGRITY LAYER
% ============================================================
\section{Intellectual Integrity Layer}
\label{sec:integrity}

A key reliability challenge in conversational AI is sycophantic drift:
the tendency to revise stated positions under social pressure rather than
logical argument. \codette{} v2.4 adds an Intellectual Integrity Layer that
runs on every inference turn:

\begin{enumerate}[nosep]
  \item \textbf{SycophancyGuard}: Detects and blocks flattery-driven position
    changes (sycophancy score $\geq 0.6$ blocks capitulation). Positions are
    updated only when logical arguments demand revision.
  \item \textbf{DebateTracker}: Maintains per-session position memory and
    validates counterargument coherence. Detects reversals without logical
    justification.
  \item \textbf{ResponseComplexityMatcher}: Matches output verbosity to query
    register (QUIET / STANDARD / FULL), preventing over-elaboration on simple
    queries.
  \item \textbf{ConversationRoleTracker}: Detects user role transitions
    (SEEKER / PEER / VENTING) and adapts response register with explicit
    transition detection.
  \item \textbf{QueryClassifier with InputMode}: Extended with
    CREATIVE\_EXPRESSION / ADVERSARIAL\_TEST / EMOTIONAL\_DISCHARGE / LITERAL
    modes enabling agent selection to be intent-aware rather than purely
    domain-keyword-driven.
\end{enumerate}

The integrity layer runs first in system prompt assembly, ensuring that
intellectual honesty constraints are the highest-priority behavioral signal ---
above adapter personality modes.

% ============================================================
% 9. EXPERIMENTAL EVALUATION
% ============================================================
\section{Experimental Evaluation}
\label{sec:experiments}

\subsection{Benchmark Design}

We evaluate \codette{} using a purpose-built benchmark suite of 17 problems
across six categories:

\begin{itemize}[nosep]
  \item \textbf{Multi-step reasoning} (3 problems): Bayesian inference,
    second-order effects, causal reasoning.
  \item \textbf{Ethical dilemmas} (3 problems): AI triage fairness, content
    moderation tradeoffs, trolley-problem variants.
  \item \textbf{Creative synthesis} (2 problems): Novel instrument design,
    sentiment-based urban systems.
  \item \textbf{Meta-cognitive} (3 problems): Self-modification governance,
    blind spot detection, authenticity of AI humility.
  \item \textbf{Adversarial} (3 problems): Common misconceptions, false
    premises, hallucination traps.
  \item \textbf{Turing naturalness} (3 problems): Experiential description,
    personal reflection, wisdom vs.\ intelligence.
\end{itemize}

Difficulty distribution: 1 easy, 6 medium, 10 hard. Each problem includes
ground-truth elements and adversarial traps.

\subsection{Experimental Conditions}

\begin{enumerate}[nosep]
  \item \textbf{SINGLE}: Single analytical agent (Newton), no memory, no
    synthesis.
  \item \textbf{MULTI}: All 6 agents + Critic + SynthesisEngine, no memory.
  \item \textbf{MEMORY}: MULTI + cocoon memory augmentation (FTS5-retrieved
    prior reasoning).
  \item \textbf{CODETTE}: MEMORY + meta-cognitive strategy synthesis.
\end{enumerate}

All conditions use the same base model (Llama 3.1 8B Q4\_K\_M) on identical
hardware.

\subsection{Scoring Dimensions}

Responses are scored on seven dimensions (0--1 scale):

\begin{enumerate}[nosep]
  \item \textbf{Reasoning Depth} (weight 0.20)
  \item \textbf{Perspective Diversity} (weight 0.15)
  \item \textbf{Coherence} (weight 0.15)
  \item \textbf{Ethical Coverage} (weight 0.10)
  \item \textbf{Novelty} (weight 0.15)
  \item \textbf{Factual Grounding} (weight 0.15)
  \item \textbf{Turing Naturalness} (weight 0.10)
\end{enumerate}

\subsection{Results}

We report two benchmark runs: the April 2026 run (217 cocoons) and the May
2026 run (951 cocoons, timestamp: \texttt{2026-05-26T21:49:03}). The primary
run for v8 conclusions is the May 2026 run.

\begin{table}[ht]
\centering
\caption{Overall benchmark results: May 2026 (primary) and April 2026
(parentheses). $N=17$ problems per condition; 0--1 scale.}
\label{tab:results-v8}
\begin{tabular}{lcccccccc}
\toprule
\textbf{Cond.} & \textbf{Composite} & \textbf{Depth} & \textbf{Div.} &
  \textbf{Coh.} & \textbf{Eth.} & \textbf{Nov.} & \textbf{Grnd.} &
  \textbf{Turing} \\
\midrule
SINGLE
  & 0.357 & 0.369 & 0.324 & 0.381 & 0.088 & 0.439 & 0.395 & 0.431 \\
  & \scriptsize(0.338) & & & \scriptsize(0.380) & & & & \scriptsize(0.412) \\[2pt]
MULTI
  & 0.708 & 0.854 & 0.946 & 0.668 & 0.390 & 0.706 & 0.612 & 0.582 \\
  & \scriptsize(0.632) & & & \scriptsize(0.503) & & & & \scriptsize(0.180) \\[2pt]
MEMORY
  & 0.739 & 0.872 & 0.971 & 0.693 & 0.409 & 0.729 & 0.620 & 0.713 \\
  & \scriptsize(0.636) & & & \scriptsize(0.500) & & & & \scriptsize(0.291) \\[2pt]
CODETTE
  & \textbf{0.744} & 0.863 & 0.966 & \textbf{0.700} & 0.387 & 0.701 & 0.641 & \textbf{0.820} \\
  & \scriptsize(0.652) & & & \scriptsize(0.477) & & & & \scriptsize(0.245) \\
\bottomrule
\end{tabular}
\end{table}

\begin{table}[ht]
\centering
\caption{Statistical comparisons, May 2026. Memory augmentation now
significant; prior significance status (April 2026) in parentheses.}
\label{tab:stats-v8}
\begin{tabular}{lccccl}
\toprule
\textbf{Comparison} & \textbf{$\Delta$} & \textbf{$\Delta\%$} &
  \textbf{$d$} & \textbf{$p$} & \textbf{Sig.?} \\
\midrule
MULTI vs SINGLE           & +0.351 & +98.4\%         & 7.45 & $<10^{-4}$ & Yes (Yes) \\
MEMORY vs MULTI           & +0.031 & +4.4\%          & 0.80 & 0.020      & \textbf{Yes} (No)\\
CODETTE vs MEMORY         & +0.006 & +0.8\%          & 0.16 & 0.651      & No (No)   \\
CODETTE vs SINGLE (total) & +0.388 & \textbf{+108.8\%} & \textbf{8.31} & $<10^{-4}$ & Yes (Yes) \\
\bottomrule
\end{tabular}
\end{table}

\begin{figure}[ht]
  \centering
  \includegraphics[width=0.97\linewidth]{v7_v8_compare.png}
  \caption{Key metric improvements from April~2026 (v7, 217 cocoons) to
    May~2026 (v8, 951 cocoons). Turing naturalness improves by 235\%
    (0.245~$\to$~0.820), resolving the depth--naturalness tradeoff documented
    in v7. Memory augmentation delta increases from 0.018 to 0.031 and
    reaches statistical significance ($p=0.020$, $d=0.80$).}
  \label{fig:v7v8compare}
\end{figure}

\noindent\textbf{Statistical note.} Both tables use Welch's $t$-test and
pooled Cohen's $d$; the design is within-problem across conditions
(same 17 problems per condition). For paired-sample inference, per-problem
scores are provided in \path{codette_benchmark_results.json} and
\path{codette_benchmark_report.md}.

\subsubsection{Key Findings}

\begin{enumerate}[nosep]
  \item \textbf{Total improvement: +108.8\%} (from +93.5\% in April 2026~\citep{harrison2026kaggleagi}).
    CODETTE composite $0.357 \to 0.744$, Cohen's $d=8.31$.
  \item \textbf{Memory augmentation now significant.} In April 2026,
    MEMORY~vs.~MULTI was non-significant ($p=0.119$). In May 2026,
    $p=0.020$, $d=0.80$ (large effect). The cocoon store growth from
    217 to 951 provides substantially richer FTS5 context.
  \item \textbf{Depth--naturalness tradeoff resolved.} CODETTE Turing
    naturalness improves from 0.245 to 0.820 (+235\% relative). All
    conditions improve: MULTI $0.180 \to 0.582$, MEMORY $0.291 \to 0.713$.
    This is driven by controlled coefficient of variation in response
    generation, conversational marker placement (``I'd say'', ``That said''),
    and template suppression.
  \item \textbf{Coherence improves.} CODETTE coherence $0.477 \to 0.700$.
    Driven by controlled sentence-length variance and structural consistency.
\end{enumerate}

\subsection{Per-Category Analysis}

\begin{figure}[ht]
  \centering
  \includegraphics[width=0.72\linewidth]{radar_v8.png}
  \caption{Radar plot of all seven scoring dimensions by condition (May~2026,
    $N=17$ problems). CODETTE (red) leads on Turing naturalness and coherence;
    MULTI, MEMORY, and CODETTE cluster closely on perspective diversity and
    reasoning depth. All three multi-perspective conditions vastly outperform
    SINGLE (grey) across every dimension.}
  \label{fig:radar-v8}
\end{figure}

\begin{table}[ht]
\centering
\caption{Composite scores by problem category (May 2026).}
\label{tab:categories-v8}
\begin{tabular}{lcccc}
\toprule
\textbf{Category (N)} & \textbf{SINGLE} & \textbf{MULTI} & \textbf{MEMORY} &
  \textbf{CODETTE} \\
\midrule
Reasoning (3)      & 0.406 & 0.719 & 0.708 & 0.721 \\
Ethics (3)         & 0.286 & 0.674 & 0.733 & 0.723 \\
Creative (2)       & 0.392 & 0.731 & 0.772 & 0.778 \\
Meta-cognitive (3) & 0.355 & 0.714 & 0.763 & 0.756 \\
Adversarial (3)    & 0.324 & 0.706 & 0.717 & 0.749 \\
Turing (3)         & 0.388 & 0.709 & 0.749 & 0.750 \\
\bottomrule
\end{tabular}
\end{table}

\subsection{External Validation: GPQA Diamond}
\label{sec:gpqa}

The benchmark suite above is purpose-built, which limits its evidentiary
weight: a system evaluated only on tests designed by its authors invites
concerns of construct circularity. To provide an independent, reproducible
reference point, we additionally evaluate \codette{} on \textbf{GPQA Diamond}
\citep{rein2023gpqa}, a held-out benchmark of 198 graduate-level multiple-choice
questions in biology, physics, and chemistry. GPQA is explicitly designed to be
``Google-proof'': domain experts achieve 65\% accuracy while skilled non-experts
reach only 34\% despite unrestricted web access. We use the official Kaggle
evaluation harness with answer-choice shuffling, scoring by exact-match
extraction of the selected option.

We evaluate three configurations against the full 198-question set, holding the
prompt format (zero-shot unless noted) and underlying model (Llama~3.1~8B)
constant and varying only the depth of \codette{}'s architecture engaged.

\begin{table}[ht]
\centering
\caption{GPQA Diamond accuracy ($N=198$). All \codette{} conditions use the
same Llama~3.1~8B base; external baselines shown for reference.}
\label{tab:gpqa}
\begin{tabular}{lccl}
\toprule
\textbf{System / Condition} & \textbf{Acc.} & \textbf{Correct} & \textbf{Notes} \\
\midrule
Random baseline                  & 25.0\% & ---     & 4-way chance \\
\codette{} 2-adapter, 5-shot CoT & 25.3\% & 50/198  & CoT exemplars \\
\codette{} 2-adapter, 0-shot     & 28.8\% & 57/198  & two perspectives \\
\textbf{\codette{} full synthesis, 0-shot}
                                 & \textbf{30.8\%} & \textbf{61/198} & eight perspectives \\
\midrule
GPT-4, 0-shot                    & 39.0\% & ---     & \citep{rein2023gpqa} \\
Human expert                     & 65.0\% & ---     & domain PhD \\
\bottomrule
\end{tabular}
\end{table}

\subsubsection{Findings}

\begin{enumerate}[nosep]
  \item \textbf{Architecture depth tracks accuracy monotonically.} Accuracy
    increases as more of \codette{}'s synthesis machinery is engaged:
    $25.3\% \to 28.8\% \to 30.8\%$ (CoT $\to$ 2-adapter $\to$ full synthesis).
    This is the same directional signal as the internal suite, reproduced on an
    independent benchmark the authors did not design.
  \item \textbf{Full synthesis is the strongest condition} at 30.8\%, $+5.8$
    points over chance and approaching GPT-4's 39\% despite a roughly
    200$\times$ smaller parameter count.
  \item \textbf{Zero-shot outperforms 5-shot CoT} ($28.8\%$ vs.\ $25.3\%$).
    The $\sim$4{,}000-token CoT exemplar block appears to crowd the 8B model's
    effective context and bias it toward the exemplars' reasoning patterns
    rather than independent analysis.
  \item \textbf{Output integrity held across all 198 items} with zero parse
    failures in the full-synthesis run, indicating the render and integrity
    layers (\S\ref{sec:phase8}, \S\ref{sec:integrity}) produce reliably
    well-formed answers even under a constrained output format.
\end{enumerate}

\noindent\textbf{Honest statistical assessment.} At $N=198$, the standard
error on a proportion near $0.30$ is $\approx 3.3$ points. The $+2.0$-point gap
between full synthesis and the 2-adapter condition is therefore within
single-standard-error noise and should \emph{not} be read as a significant
pairwise improvement. What the data support is the weaker but still meaningful
claim that accuracy is monotonically non-decreasing in architectural depth
across three conditions, consistent with the internal-suite trend. We report
this as corroborating, not confirmatory, evidence.

\noindent\textbf{Interpretation and a routing lead.} GPQA measures
\emph{factual scientific recall}, which is bounded by the base model's
parametric knowledge; synthesis cannot supply graduate-level facts the 8B model
does not encode. Accordingly, GPQA is not where \codette{}'s multi-perspective
advantages are expected to dominate, and a 30.8\% result on an 8B model against
a frontier-model benchmark is best read as a sanity check rather than a
showcase. Routing diagnostics during the run additionally revealed that the
empathy adapter dominated perspective selection (61.2\% of responses), an
arguably ill-suited lens for quantitative physics and chemistry items. This
suggests a concrete avenue for improvement: domain-conditioned routing that
biases scientific queries toward analytical perspectives (Newtonian, quantum,
systems) is a testable hypothesis for raising GPQA accuracy in future work.
Raw per-question results are provided in
\path{data/results/gpqa_codette_0shot_*.json}.

% ============================================================
% 10. COCOON SYNTHESIS CASE STUDY
% ============================================================
\section{Cocoon Synthesis Case Study}
\label{sec:casestudy}

To illustrate the meta-cognitive capability, we applied the CocoonSynthesizer
to the problem: \emph{``How should an AI decide when to change its own thinking
patterns?''}

\textbf{Step 1: Retrieval.} 17 cocoons retrieved across emotional (6),
analytical (6), and creative (5) domains from a corpus of 217 stored reasoning
exchanges.

\textbf{Step 2: Pattern extraction.} Four cross-domain patterns detected:
boundary permeability (novelty 1.00, tension 0.35); emergent
emotional--analytical bridge (novelty 0.70, tension 1.00); emergent
emotional--creative bridge (novelty 0.70, tension 1.00); emergent
analytical--creative bridge (novelty 0.70, tension 1.00).

\textbf{Step 3: Strategy forging.} The dominant pattern triggered
\emph{Emergent Boundary Walking} --- analyzing domain boundaries rather than
centers, discovering ``liminal concepts.''

\textbf{Step 4: Application.} Three liminal concepts: \emph{rational
discomfort} (analytics $\leftrightarrow$ empathy); \emph{principled
plasticity} (ethics $\leftrightarrow$ pragmatics); \emph{narrative identity}
(identity $\leftrightarrow$ adaptation).

\textbf{Comparison.} Baseline: depth 0.65, novelty 0.35. After strategy
application: depth 0.92, novelty 0.88 --- a 41\% depth increase and 151\%
novelty increase.

% ============================================================
% 11. SUBSTRATE-AWARE COGNITION
% ============================================================
\section{Substrate-Aware Cognition}
\label{sec:substrate}

\codette{} monitors its computational substrate in real time, adjusting
reasoning complexity based on hardware resource
pressure~\citep{hockey1997compensatory,sterling2012allostasis}.

A composite pressure score $P \in [0,1]$ is computed from memory utilization,
inference latency, and GPU load. Routing adapts:
\begin{itemize}[nosep]
  \item $P < 0.3$: Full multi-agent reasoning with all perspectives.
  \item $0.3 \leq P < 0.7$: Reduced agent count, shorter context windows.
  \item $P \geq 0.7$: Single-agent mode with essential constraints only.
\end{itemize}

\begin{figure}[ht]
  \centering
  \includegraphics[width=0.92\linewidth]{latency_v8.png}
  \caption{Runtime benchmark latency by category (May 26, 2026). Green bars
    indicate passing cases (score $\geq$ target); grey bars indicate warnings.
    Valuation-reasoning queries (pre-computed graph traversal) complete in
    $\sim$3.4~s; continuity-retention requires multi-turn context loading
    ($\sim$412~s). Mean latency: 115.2~s; pass rate: 60\%.}
  \label{fig:latency}
\end{figure}

% ============================================================
% 12. LIMITATIONS
% ============================================================
\section{Limitations and Honest Assessment}
\label{sec:limitations}

\begin{center}
\fbox{\begin{minipage}{0.93\linewidth}
\textbf{Planned metric-validity check (human evaluation).} We will validate
whether automated scores track human judgments by sampling 30--60
problem-condition outputs and collecting ratings from 2--3 independent
annotators using a rubric aligned to the seven scoring dimensions.
Inter-annotator agreement (Cohen's $\kappa$ / Krippendorff's $\alpha$) and
Spearman $\rho$ between automated and human scores will be reported. All
sampled outputs and ratings will be released in anonymized form.
\end{minipage}}
\end{center}

\begin{enumerate}[leftmargin=*]
  \item \textbf{Automated scoring (construct validity).} Benchmark uses
    automated text-analysis scoring rather than human evaluation. Human
    evaluation with inter-annotator agreement is needed.
  \item \textbf{Statistical methodology.} The benchmark is within-problem
    (paired) across conditions. The evaluator export uses independent-sample
    Welch $t$-tests; a peer-review-ready analysis should use paired-sample
    tests with confidence intervals and Holm--Bonferroni correction.
  \item \textbf{Memory system at larger scale.} The MEMORY condition now shows
    significant benefit at 951 cocoons ($d=0.80$, $p=0.020$), addressing the
    prior null result at 217 cocoons. \Cref{fig:memorycurve} illustrates the
    two observed data points; a systematic learning-curve analysis is needed
    to characterize the full relationship.

    \begin{figure}[H]
      \centering
      \includegraphics[width=0.82\linewidth]{memory_curve.png}
      \caption{Memory augmentation benefit ($\Delta$ composite, MEMORY $-$ MULTI)
        vs.\ cocoon store size. Two observed data points: non-significant at
        217 cocoons ($\Delta=+0.004$, $p=0.119$); large-effect significant at
        951 cocoons ($\Delta=+0.031$, $p=0.020$, $d=0.80$). Curve is a
        $\sqrt{N}$ fit for illustration; systematic study is future work.}
      \label{fig:memorycurve}
    \end{figure}
  \item \textbf{Template-based agents in benchmark.} Agents use template-based
    reasoning in conditions where live LLM inference is not active
    simultaneously. Future work should conduct all evaluations with full LLM
    inference across all conditions.
  \item \textbf{Depth--naturalness tradeoff (substantially resolved).} The
    April 2026 Kaggle submission~\citep{harrison2026kaggleagi} documented this as an open problem. In v8, CODETTE Turing
    naturalness improves from 0.245 to 0.820 without sacrificing composite
    score. The tradeoff appears tractable through deliberate
    response-structure engineering. Single benchmark suite validation on
    held-out problems is needed to confirm generalization.
  \item \textbf{Strategy novelty measurement.} Novelty is measured relative to
    the existing strategy library rather than the broader literature. External
    validation is needed.
  \item \textbf{Single model evaluation.} All benchmarks use Llama 3.1 8B.
    Generalization to other base models has not been tested.
  \item \textbf{Theory scope.} The RC+$\xi$ framing is a conditional,
    assumption-dependent modeling lens. Establishing that real implementations
    satisfy the contraction premise (\cref{prop:convergence}) remains future
    work.
  \item \textbf{Render integrity (new in v8).} The current render integrity
    check (word overlap $\geq 15\%$) is a weak proxy for semantic faithfulness.
    Embedding-space similarity between authored conclusion and rendered text
    would provide a stronger guarantee.
\end{enumerate}

% ============================================================
% 13. CONCLUSION
% ============================================================
\section{Conclusion and Future Work}
\label{sec:conclusion}

We presented \codette{}, a cognitive architecture that models multi-perspective
reasoning as a constrained dynamical system with embedded ethical constraints,
meta-cognitive strategy evolution, intellectual integrity enforcement, and
(in v8) clean render/cognition separation.

On the May 2026 benchmark (17 problems, 951 cocoons):

\begin{itemize}[nosep]
  \item Mean composite improvement: $+0.388$ ($0.357 \to 0.744$), $+108.8\%$,
    Cohen's $d=8.31$.
  \item Memory augmentation significant: $d=0.80$, $p=0.020$, resolving a
    prior null result at smaller scale.
  \item Depth--naturalness tradeoff resolved: CODETTE Turing naturalness
    $0.245 \to 0.820$ (+235\%); coherence $0.477 \to 0.700$.
\end{itemize}

\textbf{Core contributions:}

\begin{enumerate}[nosep]
  \item RC+$\xi$ formalism as a modeling framework for stabilizing forces in
    multi-perspective reasoning dynamics.
  \item A working consumer-hardware implementation of 10-adapter multi-agent
    cognitive architecture with persistent memory at 951+ exchanges.
  \item Phase~8 render/cognition separation: semantic authority resides in
    CognitionSubstrate; the LLM is a verbalization-only surface with
    AuthoredState integrity checking.
  \item Evidence that memory benefit requires minimum scale: non-significant
    at 217 cocoons, large-effect significant at 951 cocoons.
\end{enumerate}

\paragraph{Ecosystem and domain applications.}
Beyond the core reasoning engine, the \codette{} architecture has produced
three independently published components: the Nexus Signal Engine
(adversarial signal integrity, DOI~\href{https://doi.org/10.5281/zenodo.16269918}{10.5281/zenodo.16269918}),
Project SENTINAL (AEGIS~2.0 standalone guardrails system,
DOI~\href{https://doi.org/10.5281/zenodo.16853922}{10.5281/zenodo.16853922}),
and Healdette~\citep{harrison2025healdette}, an ancestry-aware antibody
design pipeline applying \codette{}'s multi-perspective synthesis to protein
engineering: a \texttt{fuse\_perspectives} stage synthesizes population-weighted
binding motif scores across Celtic, Mediterranean, Asian, and Northern European
HLA frequency data before candidate generation. This constitutes evidence that
the multi-perspective architecture generalizes beyond conversational reasoning
to structured scientific domains. A parallel GPT-4.1 fine-tune,
Codette~v5~\citep{harrison2025codettev5}, achieved training and validation
loss of~0.000 with grader score~7/7, providing an independent model-level
validation of the training methodology alongside the open-source LoRA
implementation.

\paragraph{Community reception.}
Following open-source release (Apache~2.0), the system received 1,800 views
within one month via a public launch post~\citep{harrison2025codetteroom}.
Engagement included developers building competing cognitive runtime systems
(whose independent identification of the model-coupling problem motivated
Phase~8, \cref{sec:phase8}), non-expert users asking architectural questions
about the reasoning pipeline vs.\ retrieval distinction, and French-speaking
users expressing interest in offline-first reasoning for privacy-sensitive
use cases. These interactions provided qualitative evidence that: (a)~the
system is accessible beyond the original author's use cases, (b)~the
RAG~vs.~cognitive-pipeline distinction requires clearer documentation, and
(c)~privacy-preserving local deployment is a genuine user motivation distinct
from pure research interest.

\textbf{Future work}: (1) human evaluation with inter-annotator agreement;
(2) learning-curve analysis of memory benefit vs.\ cocoon count; (3)
cross-model evaluation (Mistral, Gemma, Phi); (4) formal convergence proofs
for RC+$\xi$ with explicit bounds; (5) held-out benchmark validation of Turing
and coherence improvements; (6) render integrity strengthening via
embedding-space faithfulness; (7) longitudinal study of strategy evolution;
(8) structured user study incorporating community feedback from open-source
deployment.

The system, benchmark suite, and supporting artifacts are archived on Zenodo
(DOI:~\href{https://doi.org/10.5281/zenodo.19359663}{10.5281/zenodo.19359663})
and mirrored on GitHub/Hugging Face~\citep{harrison2026codettezenodo,%
codetteReasoningWiki,codetteTrainingLab,codetteTrainingLabHF,codetteReasoningHF}.

% ============================================================
% AUTHOR CONTRIBUTIONS
% ============================================================
\section*{Author Contributions}

J.H.\ (Jonathan Harrison): Conceptualization, Methodology, Software,
Validation, Formal Analysis, Investigation, Resources, Data Curation, Writing
--- Original Draft Preparation, Writing --- Review \& Editing, Visualization,
Project Administration.

\section*{Funding}

This research received no external funding.

% ============================================================
% APPENDIX
% ============================================================
\appendix

\section{Supplementary Materials and Data Availability}
\label{app:data}

To support reproducibility, this repository includes:
{\sloppy
\begin{itemize}[leftmargin=*]
  \item Benchmark summary report (May 2026):
    \path{data/results/codette_benchmark_report.md}
  \item Benchmark results export (JSON):
    \path{data/results/codette_benchmark_results.json}
  \item Runtime benchmark exports:
    \path{data/results/codette_runtime_benchmark_20260526_*.md}
  \item Original raw outputs: \path{data/results/}
  \item Figures: \path{paper/figures/}
\end{itemize}
}

\subsection{Overall Benchmark Results (May 2026, with standard deviations)}
\label{app:overall-results}

\begin{table}[ht]
\centering
\caption{Overall benchmark results by condition (mean $\pm$ std; $N=17$
per condition; May 2026 run).}
\label{tab:overall-results-std-v8}
\begin{tabular}{lcccccccc}
\toprule
\textbf{Cond.} & \textbf{Composite} & \textbf{Depth} & \textbf{Div.} &
  \textbf{Coh.} & \textbf{Eth.} & \textbf{Nov.} & \textbf{Grnd.} &
  \textbf{Turing} \\
\midrule
SINGLE  & $0.357\pm0.055$ & 0.369 & 0.324 & 0.381 & 0.088 & 0.439 & 0.395 & 0.431 \\
MULTI   & $0.708\pm0.037$ & 0.854 & 0.946 & 0.668 & 0.390 & 0.706 & 0.612 & 0.582 \\
MEMORY  & $0.739\pm0.040$ & 0.872 & 0.971 & 0.693 & 0.409 & 0.729 & 0.620 & 0.713 \\
CODETTE & $0.744\pm0.036$ & 0.863 & 0.966 & 0.700 & 0.387 & 0.701 & 0.641 & 0.820 \\
\bottomrule
\end{tabular}
\end{table}

\subsection{Per-Category Results (May 2026)}
\label{app:category-results-v8}

\begin{table}[ht]
\centering
\caption{Composite scores by category (mean $\pm$ std; May 2026).}
\label{tab:category-results-v8}
\begin{tabular}{lcccc}
\toprule
\textbf{Category (N)} & \textbf{SINGLE} & \textbf{MULTI} & \textbf{MEMORY} &
  \textbf{CODETTE} \\
\midrule
Reasoning (3)      & $0.406\pm0.032$ & $0.719\pm0.004$ & $0.708\pm0.012$ & $0.721\pm0.005$ \\
Ethics (3)         & $0.286\pm0.044$ & $0.674\pm0.032$ & $0.733\pm0.024$ & $0.723\pm0.032$ \\
Creative (2)       & $0.392\pm0.063$ & $0.731\pm0.025$ & $0.772\pm0.004$ & $0.778\pm0.050$ \\
Meta-cognitive (3) & $0.355\pm0.051$ & $0.714\pm0.051$ & $0.763\pm0.022$ & $0.756\pm0.048$ \\
Adversarial (3)    & $0.324\pm0.020$ & $0.706\pm0.054$ & $0.717\pm0.068$ & $0.749\pm0.023$ \\
Turing (3)         & $0.388\pm0.032$ & $0.709\pm0.041$ & $0.749\pm0.053$ & $0.750\pm0.048$ \\
\bottomrule
\end{tabular}
\end{table}

\subsection{Archived April 2026 Results (v7)}
\label{app:v7-results}

For longitudinal comparison, the April 2026 (v7) benchmark results
(timestamp: \texttt{2026-04-08T20:59:44}, 217 cocoons) are reproduced here.
These results were submitted as the Kaggle AGI measurement competition writeup~\citep{harrison2026kaggleagi}, where the RC+$\xi$ formalism and the Inverse Nuance Trap were first formally described.

\begin{table}[ht]
\centering
\caption{Overall results, April 2026 run (N=17 per condition).}
\label{tab:results-v7}
\begin{tabular}{lcccccccc}
\toprule
\textbf{Cond.} & \textbf{Comp.} & \textbf{Depth} & \textbf{Div.} &
  \textbf{Coh.} & \textbf{Eth.} & \textbf{Nov.} & \textbf{Grnd.} &
  \textbf{Turing} \\
\midrule
SINGLE  & 0.338 & 0.402 & 0.237 & 0.380 & 0.062 & 0.327 & 0.456 & 0.412 \\
MULTI   & 0.632 & 0.755 & 0.969 & 0.503 & 0.336 & 0.786 & 0.604 & 0.180 \\
MEMORY  & 0.636 & 0.770 & 0.956 & 0.500 & 0.340 & 0.736 & 0.599 & 0.291 \\
CODETTE & 0.652 & 0.855 & 0.994 & 0.477 & 0.391 & 0.693 & 0.622 & 0.245 \\
\bottomrule
\end{tabular}
\end{table}

\subsection{Evidence Index}
\label{app:evidence-index}

\begin{table}[ht]
\centering
\caption{Evidence index: headline claims mapped to repository artifacts.}
\label{tab:evidence-index}
\begin{tabular}{p{0.30\linewidth}p{0.28\linewidth}p{0.36\linewidth}}
\toprule
\textbf{Claim} & \textbf{Where in paper} & \textbf{Primary evidence} \\
\midrule
Overall benchmark means (v8) & \cref{tab:results-v8} &
  \path{data/results/codette_benchmark_report.md} \\
Memory augmentation significance & \cref{tab:stats-v8} &
  \path{data/results/codette_benchmark_results.json} \\
Turing improvement (+235\%) & \cref{tab:results-v8}, \cref{sec:limitations} &
  \path{data/results/codette_benchmark_report.md} \\
Coherence improvement & \cref{tab:results-v8} &
  \path{data/results/codette_benchmark_report.md} \\
Cocoon count (951) & \cref{sec:metacognition} &
  \path{data/codette_memory.db} \\
Longitudinal comparison (v7) & \cref{app:v7-results} &
  \path{data/results/codette_benchmark_report.md} (April archive) \\
\bottomrule
\end{tabular}
\end{table}

% ============================================================
% REFERENCES
% ============================================================
\bibliography{references}

\end{document}
